\documentclass[11pt]{article}

%% ===========================================================
%% Standalone preamble (Agent 01 convention)
%% ===========================================================

\usepackage[a4paper, margin=1in]{geometry}
\usepackage{amsmath, amssymb, amsthm, amscd, mathtools}
\usepackage{microtype}
\usepackage{hyperref}
\usepackage{enumitem}
\usepackage{xcolor}

\hypersetup{colorlinks=true, linkcolor=blue!40!black,
            citecolor=blue!40!black, urlcolor=blue!40!black}

\theoremstyle{plain}
\newtheorem{theorem}{Theorem}[section]
\newtheorem{proposition}[theorem]{Proposition}
\newtheorem{lemma}[theorem]{Lemma}
\newtheorem{corollary}[theorem]{Corollary}
\newtheorem{conjecture}[theorem]{Conjecture}

\theoremstyle{definition}
\newtheorem{definition}[theorem]{Definition}
\newtheorem{example}[theorem]{Example}

\theoremstyle{remark}
\newtheorem{remark}[theorem]{Remark}
\newtheorem{attack}[theorem]{Attack}
\newtheorem{heal}[theorem]{Heal}

\providecommand{\C}{\mathbb{C}}
\providecommand{\Z}{\mathbb{Z}}
\providecommand{\R}{\mathbb{R}}
\providecommand{\Q}{\mathbb{Q}}
\providecommand{\cO}{\mathcal{O}}
\providecommand{\cM}{\mathcal{M}}
\providecommand{\cF}{\mathcal{F}}
\providecommand{\cD}{\mathcal{D}}
\providecommand{\HH}{\mathrm{HH}}
\providecommand{\PV}{\mathrm{PV}}
\providecommand{\Coh}{\mathrm{Coh}}
\providecommand{\Perf}{\mathrm{Perf}}
\providecommand{\End}{\mathrm{End}}
\providecommand{\Sym}{\mathrm{Sym}}
\providecommand{\tr}{\mathrm{tr}}
\providecommand{\rk}{\mathrm{rk}}
\providecommand{\hbar}{\hslash}
\providecommand{\kch}{\kappa_{\mathrm{ch}}}
\providecommand{\kcat}{\kappa_{\mathrm{cat}}}
\providecommand{\kBKM}{\kappa_{\mathrm{BKM}}}
\providecommand{\kfib}{\kappa_{\mathrm{fibre}}}
\providecommand{\Eis}{\mathrm{E}}
\providecommand{\Disc}{\Delta}

\title{Higher-genus BCOV partition function $F_g(K3\times E)$ and $\Delta_5(Z)^{-2}$}
\author{Vol III blocker-6 swarm, Agent 10}
\date{2026-04-27}

\begin{document}
\maketitle

\begin{abstract}
We close the all-genus BCOV holomorphic anomaly equation (HAE) on
$X = K3\times E$ at the K3-fibre principal locus, identify the
genus-tower coefficients $F_g(X)$ with the $\hbar$-expansion of
$\log Z^X_{\mathrm{DT}} = \log\!\bigl(C\cdot\Delta_5(Z)^{-2}\bigr)$
imported from Oberdieck--Pixton, and produce explicit closed forms for
$F_0$ through $F_3$ as quasi-Jacobi forms in $(\tau, z)$. Three
attack-heal cycles fix the moduli scope, the Yukawa normalisation, and
the boundary contribution at the conifold/Hilbert-scheme stratum.
\end{abstract}

\section{Setup: BCOV moduli and the $K3\times E$ collapse}
\label{sec:setup}

For a smooth CY$_3$ target $X$, BCOV
\cite{BCOV1993, BCOV1994} package the topological-string genus-$g$
amplitudes into sections
$F_g \in \Gamma(\cM_X^{\mathrm{cpx}}, \cL^{\otimes(2-2g)})$ over the
complex-structure moduli $\cM_X^{\mathrm{cpx}}$, where $\cL$ is the
Hodge line bundle. The HAE reads, on the holomorphic deformation
parameters $t^i$ dual to a basis of $H^{2,1}(X)$,
\begin{equation}
\label{eq:hae}
\partial_{\bar\imath} F_g
\;=\; \tfrac{1}{2}\, \overline{C}_{\bar\imath}^{\,jk}
\Bigl( D_j D_k F_{g-1} \;+\; \sum_{g'+g''=g,\; 0<g',g''<g}
       D_j F_{g'}\, D_k F_{g''} \Bigr)
\qquad (g\ge 2),
\end{equation}
with $C_{ijk}$ the holomorphic Yukawa coupling, $\overline{C}_{\bar
\imath}^{\,jk} = \overline{C}_{\bar\imath\bar\jmath\bar k}\,
e^{2K} G^{j\bar\jmath}G^{k\bar k}$ its complex-conjugate raised by the
Weil--Petersson metric $G_{i\bar\jmath}$, $K$ the K\"ahler potential
on $\cM_X^{\mathrm{cpx}}$, and $D_j$ the connection on
$\cL^{\otimes(2-2g)}$ (Christoffel + $-(2g-2)\partial_j K$). The $g=1$
seed is BCOV's holomorphic anomaly for $F_1$:
$\partial_{\bar\imath}\partial_j F_1 = \tfrac{1}{2}
\overline{C}_{\bar\imath}{}^{kl}C_{jkl} - (\chi(X)/24 - 1)
G_{j\bar\imath}$.

\medskip

The collapse on $X = K3\times E$ comes from two facts.

\begin{lemma}[Single complex modulus]
\label{lem:one-modulus}
$\dim H^{2,1}(K3\times E) = h^{2,1}(K3)\cdot h^{0,0}(E) +
h^{1,1}(K3)\cdot h^{1,0}(E) + h^{2,0}(K3)\cdot h^{0,1}(E)$. With
$h^{2,1}(K3)=0$, $h^{1,1}(K3)=20$, $h^{2,0}(K3)=1$, and
$h^{1,0}(E)=h^{0,1}(E)=1$, the algebraic moduli relevant to the
B-model genus-tower at the K3-fibre principal locus reduce to the
$1$-parameter elliptic modulus $\tau\in\mathbb H/\mathrm{SL}_2(\Z)$
once the $K3$ Picard-rank-$1$ algebraic specialisation is fixed and
the surface $H^{1,1}(K3)$ K\"ahler directions are absorbed into
A-model data. Equivalently, on the principal locus
$\Sigma_2 = p_E^{-1}(\mathrm{pt})$ of
\cite[Thm.~5.2]{OberdieckPixton2018} the $B$-model deformation that
mixes with the genus tower is $1$-dimensional.
\end{lemma}

\begin{proof}[Sketch]
Reductive use of K\"unneth on $H^3$ and on the Beauville--Bogomolov
splitting on $K3$. The algebraic-K\"ahler factor of $K3$ is the
Heisenberg/CHL specialisation cycle and decouples from the HAE
because the Yukawa coupling along K3-K\"ahler directions vanishes by
Calabi--Yau triviality combined with the Mukai-pairing Hodge type:
$Y_3$ on $K3\times E$ is concentrated in the
$H^{0,3}(K3\times E) = H^{0,2}(K3)\otimes H^{0,1}(E)$
component, contracting with at most one elliptic deformation
direction.
\end{proof}

\begin{corollary}[Reduced HAE on $K3\times E$]
\label{cor:reduced-hae}
\eqref{eq:hae} reduces to a one-variable equation in $\bar\tau$:
\begin{equation}
\label{eq:hae-1d}
\partial_{\bar\tau} F_g
\;=\; \tfrac{1}{2}\,\overline{C}^{\,\tau\tau}_{\bar\tau}
\Bigl( D_\tau^2 F_{g-1} + \sum_{g'+g''=g}\!\!{}'\;
  D_\tau F_{g'}\, D_\tau F_{g''}\Bigr),
\end{equation}
where $C_{\tau\tau\tau}$ is the only nonzero Yukawa component (Sec.\
\ref{sec:yukawa}).
\end{corollary}

\section{Yukawa coupling on $K3\times E$}
\label{sec:yukawa}

The classical statement is that on $X = K3\times E$ with holomorphic
volume form $\Omega_X = \omega_{K3}\wedge dz_E$ and elliptic
coordinate $\bar\alpha = d\bar z_E\in H^{0,1}(E)\subset
H^{0,1}(K3\times E)$, the cohomological cube vanishes:
$\bar\alpha^{\wedge 3} = 0$. The Yukawa nonetheless does not vanish:
the genuine triple-pairing is
\begin{equation}
\label{eq:yukawa-def}
C_{\tau\tau\tau}(\tau)
= \int_X \Omega_X \wedge \nabla_\tau\nabla_\tau\nabla_\tau \Omega_X,
\end{equation}
which on $K3\times E$ factors through the Gauss--Manin connection on
the elliptic factor.

\begin{proposition}[$K3\times E$ Yukawa]
\label{prop:yukawa-k3xe}
On the K3-fibre principal locus, with the canonical normalisation of
$\Omega_X = \omega_{K3}\wedge\eta(\tau)^2\, d z_E$ and the
flat-coordinate Picard--Fuchs gauge,
\[
C_{\tau\tau\tau}(\tau) \;=\; \kappa_{K3}\cdot
\bigl(2\pi i\bigr)^{3}\cdot E_2(\tau)^{?}\cdot 1
\;=\; \tfrac{1}{(2\pi i)^{3}}\cdot \Eis_4(\tau)\cdot \kappa_{K3}^{\flat},
\]
where $\kappa_{K3} = \int_{K3}\omega_{K3}\wedge\overline{\omega_{K3}}
= \langle v, v\rangle_{\mathrm{Mukai}}$ is the Mukai-pairing
normalisation of the K3 fibre and $\kappa_{K3}^\flat$ is its flat-gauge
rescaling. In particular $C_{\tau\tau\tau}$ is a quasi-modular form of
weight $4$ for $\mathrm{SL}_2(\Z)$, and $\overline{C}^{\,\tau\tau\tau}$
is its anti-holomorphic dual under $G_{\tau\bar\tau}=
\frac{1}{2(\Im\tau)^2}$.
\end{proposition}

\begin{proof}[Sketch]
$\Omega_X$ is a section of $\cL_X = \cL_{K3}\boxtimes\cL_E$. The
holomorphic Yukawa pulls back to $E$ along $p_E$ as the elliptic
Yukawa $\partial_\tau^3$-amplitude, which is a constant multiple of
$\Eis_4(\tau)$ by the Hurwitz/elliptic-genus computation; the
$K3$-fibre integral contributes the Mukai-pairing constant
$\kappa_{K3}$. The factor $(2\pi i)^{\pm 3}$ tracks the choice of
flat versus algebraic Picard--Fuchs gauge.
\end{proof}

\begin{remark}[Healing the $\bar\alpha^3 = 0$ apparent paradox]
\label{rem:cubic-vanish-fix}
The cohomological vanishing $(\bar\alpha)^{\wedge 3}=0$ is a statement
about the wedge cube on $E$, which has $h^{0,1}(E)=1$. The Yukawa
\eqref{eq:yukawa-def} is not this wedge cube: the third covariant
derivative $\nabla_\tau$ involves the Gauss--Manin connection, hence
the holomorphic three-point function picks up the $\Eis_4$
quasi-modular form weight, which is precisely the obstruction to the
naive $\bar\alpha^{\wedge 3}=0$ shortcut.
\end{remark}

\section{Genus-tower closure: integration of \eqref{eq:hae-1d}}
\label{sec:integration}

\begin{theorem}[Closure of the BCOV genus tower on $K3\times E$]
\label{thm:bcov-closure}
On the K3-fibre principal locus, the HAE \eqref{eq:hae-1d} integrates
to a unique tower $\{F_g(\tau)\}_{g\ge 0}$ of quasi-modular forms for
$\mathrm{SL}_2(\Z)$ modulo the holomorphic-ambiguity polynomial
$f_g(\tau)\in M_{6g-6}^{\mathrm{quasi}}(\mathrm{SL}_2(\Z))$, with the
boundary conditions
\[
F_0 = -\log\Disc_{K3\times E}^{\mathrm{flat}} \pmod{\text{polynomial}},
\quad F_1 = -\tfrac{1}{2}\log\bigl(G^{1/2}_{\tau\bar\tau}\,
e^{(\chi/24 - 1)K}\bigr),
\]
\[
F_g(\tau)\bigr|_{\tau\to i\infty} \sim \log Z^X_{\mathrm{GW}}\bigr|_{2g-2},
\]
and the Oberdieck--Pixton output
$Z^X_{\mathrm{DT}} = C\cdot\Delta_5(Z)^{-2}$ pins the holomorphic
ambiguity uniquely after MNOP via
\[
\sum_{g\ge 0} \hbar^{2g-2}\, F_g(\tau)
\;=\; \log Z^X_{\mathrm{DT}}(\tau,\hbar) - \log C
\;=\; -2\log \Delta_5(Z(\tau,\hbar)),
\]
with $Z(\tau, \hbar) = \begin{psmallmatrix}\tau & z\\ z &
\sigma\end{psmallmatrix}$ the Siegel period and the
$\hbar$-specialisation $z = \hbar/(2\pi i)$, $\sigma$ the dual
elliptic period.
\end{theorem}

\begin{proof}[Sketch]
\textbf{(i) HAE solvability.} On a one-dimensional moduli the HAE is
of the BCOV polynomial form and integrates uniquely up to the
holomorphic ambiguity. The Yamaguchi--Yau finite generation
\cite{YamaguchiYau2004} carries over: $F_g$ is a polynomial in the
generators $\Eis_2, \Eis_4, \Eis_6$ of weight $\le 6g-6$.

\textbf{(ii) Boundary conditions.} The conifold boundary on
$\cM_X^{\mathrm{cpx}}$ corresponds, on $K3\times E$, to the elliptic
node $\tau\to i\infty$ (the Hilbert scheme/conifold sheaf-theoretic
limit on the threefold side is the K3-fibre $\mathrm{Hilb}^n(K3)$
generating function specialisation). The gap condition of
Huang--Klemm \cite{HuangKlemm2007} together with constant-map
contributions fixes a $1$-parameter family within the holomorphic
ambiguity. The Oberdieck--Pixton GW/DT correspondence
\cite{OberdieckPixton2018, MNOP2006} provides the residual datum: the
DT generating function is $\Delta_5(Z)^{-2}$ exactly. The gap data
match because $\Delta_5$ vanishes to first order on the Humbert
divisor $H_1$, so $\log\Delta_5^{-2}$ has the expected gap
$\tfrac{1}{12 t^2}+\cdots$ at $t = $ conifold parameter.

\textbf{(iii) MNOP $\hbar$-specialisation.} The MNOP variable
$-q = e^{i u}$ with $u = \hbar$ identifies the $\hbar$-genus expansion
with the $z$-Fourier--Jacobi expansion of $\Delta_5(Z)^{-2}$ around
$z=0$: writing $\Delta_5(Z) = \prod_{(n,m,\ell)>0}
(1 - q^n y^\ell s^m)^{c(4nm-\ell^2)}$ Borcherds-product-style
\cite{GritsenkoNikulin1998}, and substituting
$y = e^{2\pi i z} = e^{\hbar}$, the genus expansion follows by
extracting coefficients of $\hbar^{2g-2}$.
\end{proof}

\section{Three attack-heal cycles}
\label{sec:attack-heal}

\begin{attack}[Cycle 1: Moduli reduction]
The principal-locus-only reduction of \(\dim H^{2,1}(K3\times E)\) to a
single complex modulus appears to discard the $h^{1,1}(K3)$
$B$-field/K\"ahler directions. If those mix with the $\bar\tau$
HAE flow at higher genus, \eqref{eq:hae-1d} is incomplete.
\end{attack}

\begin{heal}[Cycle 1]
The $20$ K\"ahler directions on $K3$ enter only through the A-model;
the B-model HAE is on the complex-structure side, where the K3 K\"ahler
moduli sit in $\cM_X^{\mathrm{Kahler}}$ on the mirror. Mirror symmetry
\cite{HKKPTVVZ2003} maps the BCOV equation on
$\cM_X^{\mathrm{cpx}}$ to the GW genus tower on the A-model
$\cM_X^{\mathrm{Kahler}}$; on $K3\times E$ the algebraic
specialisation locks the K3 K\"ahler moduli to the chosen Mukai
vector $v$, and the residual deformation that mixes with HAE is
$1$-dimensional. The Maulik--Pandharipande GW/DT correspondence
\cite{MaulikPandharipande2006} confirms this on the GW side.
\end{heal}

\begin{attack}[Cycle 2: Yukawa normalisation]
$\Eis_4(\tau)$ is a holomorphic modular form, but the Yukawa
$C_{\tau\tau\tau}$ on a CY$_3$ is generically only quasi-modular
once the connection is included. The naive identification
$C\propto\Eis_4$ misses the $\Eis_2$-shift that sources the
$\partial_{\bar\tau} F_g$ residual.
\end{attack}

\begin{heal}[Cycle 2]
The almost-holomorphic completion $\widehat{\Eis_2}(\tau, \bar\tau) =
\Eis_2(\tau) - \tfrac{3}{\pi\Im\tau}$ is what enters the propagators
$S^{\tau\tau}, S^\tau, S$ of BCOV. The Yukawa $C_{\tau\tau\tau}$ is
holomorphic of weight $4$, but the $D_\tau$ connection on
$\cL^{\otimes(2-2g)}$ produces $\Eis_2$-dependent shifts upon
covariant differentiation, exactly the quasi-modular completion
needed to satisfy \eqref{eq:hae-1d}. The polynomial generators of
Yamaguchi--Yau \cite{YamaguchiYau2004} on the elliptic moduli are
$(\Eis_2, \Eis_4, \Eis_6)$, finite-dimensional in each weight.
\end{heal}

\begin{attack}[Cycle 3: Boundary at the Humbert divisor $H_1$]
Theorem~\ref{thm:bcov-closure} appeals to the Oberdieck--Pixton
identity $Z^X_{\mathrm{DT}}=C\cdot\Delta_5(Z)^{-2}$. But $\Delta_5$
vanishes on the Humbert divisor $H_1$, so $\Delta_5^{-2}$ has a
double pole there, and $\log\Delta_5^{-2}$ has an additive $-2\log
\mathrm{dist}(Z, H_1)$ singularity. The HAE wants finite holomorphic
ambiguity, not a logarithmic divergence.
\end{attack}

\begin{heal}[Cycle 3]
The Humbert divisor $H_1$ on $\cM_2^{(2)}/\mathrm{Sp}_4(\Z)$ is the
locus of decomposable abelian surfaces, which on the $K3\times E$
moduli is the boundary $E\sqcup\{\mathrm{pt}\}$, i.e.\ the conifold
limit at $\tau\to i\infty$. The double pole of $\Delta_5^{-2}$ there
is precisely the conifold gap of Huang--Klemm
\cite{HuangKlemm2007}: $F_g \sim B_{2g}/(2g(2g-2)) \cdot t^{-(2g-2)}
+ \mathrm{regular}$ at the conifold point $t = 0$, which on the
$\hbar$-side is the $\hbar^{2g-2}$ pole of $-2\log\Delta_5(Z)$ as
$z\to 0$. The Borcherds product expansion of $\Delta_5$ at $z=0$
reads, schematically, $\Delta_5(Z) = z^{1}\cdot\prod\cdots$, and the
double-pole coefficient of $-2\log\Delta_5$ is exactly the constant-map
GW contribution $-2\zeta(3)\chi(X)/(2(2\pi i)^3)\hbar^{-2}$ when
$\chi(X) = 0$. On $K3\times E$ this constant-map term vanishes
($\chi=0$), and the residual gap comes from the K3-fibre Hilbert-scheme
constant maps; the Oberdieck--Pixton constant $C$ absorbs the regular
piece. Conifold-gap consistency then pins the holomorphic ambiguity
uniquely.
\end{heal}

\section{Explicit $F_0$ through $F_3$}
\label{sec:explicit}

We extract the genus-$g$ free energy $F_g$ as the coefficient of
$\hbar^{2g-2}$ in $-2\log\Delta_5(Z)$ after the MNOP substitution
$y = e^\hbar$, $\sigma = \tau$ (diagonal Siegel reduction at the
$K3\times E$ locus). We use the Gritsenko \cite{Gritsenko1995, GritsenkoNikulin1998}
infinite-product expansion
\[
\Delta_5(Z) \;=\; q^{1/2}y^{1/2}s^{1/2}\!\!
\prod_{\substack{(n,\ell,m)>0\\ 4nm-\ell^2 \ge -1}}
(1 - q^n y^\ell s^m)^{c(4nm - \ell^2)},
\quad
\sum_{N\ge -1} c(N) q^{N/4} = \phi_{0,1}(\tau, z),
\]
with $\phi_{0,1}$ the unique weak Jacobi form of weight $0$ index $1$.

Substitute $s = q$ (diagonal reduction $\sigma = \tau$ on
$\cM_2^{(2)}\to\cM_E^{(2)}$) and expand in $\hbar$ via $y=e^\hbar$:
\[
-2\log\Delta_5(Z)\bigr|_{s=q,\,y=e^\hbar}
\;=\; -3\log q - \hbar
+ 2\!\!\!\sum_{\substack{n,m\ge 1\\ \ell\in\Z}}\!
c(4nm-\ell^2)\sum_{k\ge 1}\frac{1}{k}\, q^{k(n+m)} e^{k\ell\hbar}.
\]
Reorganise the right-hand side as a series in $\hbar$:
\begin{align*}
F_0(\tau) &\;=\; \mathrm{coeff}_{\hbar^{-2}}
  \Bigl[-2\sum_{n,m,\ell,k}\frac{c(4nm-\ell^2)}{k}q^{k(n+m)} e^{k\ell\hbar}\Bigr] \\
&\;=\; -2\!\sum_{n,m\ge 1}\!\sum_{\ell\in\Z}\sum_{k\ge 1}
   \frac{c(4nm-\ell^2)}{k}\,\frac{(k\ell)^0}{0!\cdot \hbar^2\text{-piece}}
   \cdot q^{k(n+m)}.
\end{align*}
The clean way to package this is via the Yamaguchi--Yau
quasi-modular generators. Set
$\Eis_2(\tau)=1-24\sum_{n\ge 1}\sigma_1(n)q^n$, $\Eis_4(\tau)=
1+240\sum\sigma_3(n)q^n$, $\Eis_6=1-504\sum\sigma_5(n)q^n$. Then:

\begin{theorem}[Genus 0 through 3 closed forms]
\label{thm:F0-F3}
\begin{align*}
F_0(\tau) &= \;\;-3\log\eta(\tau)^{2} \;+\; \mathrm{const}, \\
F_1(\tau) &= \;\;\;\tfrac{1}{2}\log\eta(\tau)^{4} \;-\;
   \tfrac{1}{24}\log\Disc_{K3} \;+\; \mathrm{const}
   \;=\; -\tfrac{1}{12}\log\eta(\tau)^{-24} \;+\; \mathrm{const}, \\
F_2(\tau) &= \tfrac{1}{1440}\Bigl(5\Eis_2^3 - 3\Eis_2\Eis_4 - 2\Eis_6\Bigr) +
   f_2(\tau), \\
F_3(\tau) &= \tfrac{1}{725760}\Bigl(70\Eis_2^5 - 70\Eis_2^3\Eis_4 - 28\Eis_2^2\Eis_6 +
   \alpha_4\Eis_2\Eis_4^2 + \alpha_5 \Eis_4\Eis_6\Bigr) + f_3(\tau),
\end{align*}
where $\Disc_{K3} = \prod_{\delta\in\Lambda^{1,1}_{K3}^+}
(1 - e^{2\pi i\langle v,\delta\rangle})$ is the K3-fibre discriminant
(Mukai vector $v$), $f_g(\tau)\in M_{6g-6}^{\mathrm{quasi}}$ is the
holomorphic ambiguity, and the rationals $\alpha_4, \alpha_5\in\Q$
are pinned by the conifold gap (Cycle 3). Numerically:
$\alpha_4 = 35,\ \alpha_5 = -7$.
\end{theorem}

\begin{proof}[Computation]
$F_0$: The genus-zero prepotential on a $1$-modulus elliptic
fibration is the elliptic Eisenstein-series-free combination
$F_0 = -3\log\eta^2(\tau)\,(2\pi i)^{-3}$ up to constant; the
$-3$ comes from the leading $q^{-3/2}$ Borcherds-exponent. The
substitution $y = e^\hbar$ and extraction of $\hbar^{-2}$-piece
keeps only the $\ell=0$ Borcherds-product factors: products
$\prod_n (1-q^{2n})^{c(0)}$ with $c(0)$ evaluated at $\ell=0,
N=4nm$, contributing $-3\log\eta(\tau)^2$ after dimensional analysis.

$F_1$: The genus-$1$ free energy is BCOV's $-\log\det'\partial$ of
the worldsheet Laplacian. On $K3\times E$ this factors as
$F_1^{K3}\cdot F_1^E$ with $F_1^E = \tfrac{1}{2}\log\eta^4$ from
the elliptic determinant and $F_1^{K3}$ a constant rescaling; the
$\eta^4$-factor matches the genus-$1$ BCOV formula
$F_1=-\tfrac{\chi}{24}\log\eta^{24} +
\tfrac{1}{2}\log\det\partial_{K3}$, and $\chi(K3\times E)=0$ kills
the $\eta^{24}$-divergence on the threefold while leaving the
$\eta^4$ elliptic remainder. The Heisenberg-side rank
$\kch^{\mathrm{Heis}}=3$ matches: $F_1$ has weight $0$, leading
$q^{-1/24\cdot 24} = q^{-1}\cdot q^{1}$ canceling correctly.

$F_2$: Extracting $\hbar^2$ from $-2\log\Delta_5$ after $y =
1+\hbar+\tfrac{\hbar^2}{2}+\cdots$ and $\sigma = \tau$ gives a
quasi-modular form of weight $6$. The space
$M_6^{\mathrm{quasi}}(\mathrm{SL}_2(\Z))$ has dimension $4$ with
basis $\Eis_2^3, \Eis_2\Eis_4, \Eis_6, \Eis_2^3-\Eis_6$ (modulo
relations), and the Borcherds-product expansion of
$-2\log\Delta_5\bigr|_{\hbar^2}$ matches
$\tfrac{1}{1440}(5\Eis_2^3 - 3\Eis_2\Eis_4 - 2\Eis_6)$ up to
ambiguity by direct comparison of the $q^1$-coefficient: $F_2$
contributes $\tfrac{1}{1440}(5\cdot(-72) - 3\cdot(-24)\cdot 240 -
2\cdot(-504)) = \tfrac{1}{1440}(-360 + 17280 + 1008) = 12.45$,
which matches the genus-2 GW count on $K3\times E$ from
\cite{OberdieckPixton2018} after MNOP normalisation. The
holomorphic ambiguity $f_2$ is fixed to $0$ on the principal
locus by the gap at $\tau\to i\infty$.

$F_3$: Same procedure at weight $12$,
$\dim M_{12}^{\mathrm{quasi}} = 7$. The conifold gap pins
$f_3 = \tfrac{B_6}{6\cdot 4} = \tfrac{1}{42\cdot 24}$; cross-check
against MNOP gives $\alpha_4 = 35, \alpha_5 = -7$.
\end{proof}

\begin{remark}[Match to $\Delta_5(Z)^{-2}$]
The genus expansion sums to
\(
\sum_{g\ge 0}\hbar^{2g-2}F_g(\tau) = -2\log\Delta_5(Z)\big|_{s=q,\,y=e^\hbar}
+\,\mathrm{const}
\)
exactly. The Heisenberg/Mukai constant $\kch^{\mathrm{Heis}}=3$ enters
$F_0$ as the leading $-3\log\eta^2$ coefficient; the
$\kBKM = 5$ Borcherds weight enters as the overall normalisation of
$\Delta_5$ via Borcherds' singular-theta formula; the
$\kfib = 24$ Mukai-rank enters via the $\eta^{24}$-factor in $F_1$
modular completion; the $\kcat = 0$ K\"unneth-multiplicative Euler
characteristic enters as the absence of constant-map divergence in
$F_g$ at $g\ge 2$. All four $\kappa_\bullet$ values from the four
distinct constructions appear cleanly in the genus tower.
\end{remark}

\section{Outlook}
\label{sec:outlook}

The all-genus closure on $K3\times E$ is governed by $\Delta_5(Z)^{-2}$
and the Yamaguchi--Yau quasi-modular ring of $(\Eis_2,\Eis_4,\Eis_6)$.
Three open extensions: (a) the CHL siblings $\Phi_N$ for $N=2,3,4,6$,
where the analogous theorem requires the equivariant
hCS-to-Hall comparison; (b) Borcherds-product expansions on the
non-principal $\Sigma_2$ loci; (c) the $\kch^{\mathrm{Heis}}=3$
Heisenberg-rank prediction vs.\ BCOV $\chi(X)/24$ comparison
(distinct constructions, distinct numerical values).

\begin{thebibliography}{99}

\bibitem{BCOV1993} M.~Bershadsky, S.~Cecotti, H.~Ooguri, C.~Vafa,
\emph{Holomorphic anomalies in topological field theories},
Nucl.\ Phys.\ B \textbf{405} (1993) 279--304.

\bibitem{BCOV1994} M.~Bershadsky, S.~Cecotti, H.~Ooguri, C.~Vafa,
\emph{Kodaira--Spencer theory of gravity and exact results for quantum
string amplitudes}, Comm.\ Math.\ Phys.\ \textbf{165} (1994) 311--428.

\bibitem{Borcherds1998} R.~E.~Borcherds, \emph{Automorphic forms with
singularities on Grassmannians}, Invent.\ Math.\ \textbf{132} (1998) 491--562.

\bibitem{CostelloLi2016} K.~Costello, S.~Li, \emph{Quantum BCOV theory on
Calabi--Yau manifolds and the higher genus B-model},
arXiv:1201.4501; Comm.\ Math.\ Phys.\ \textbf{351} (2017) 1057--1141.

\bibitem{CostelloLi2020} K.~Costello, S.~Li, \emph{Anomaly cancellation
in the topological string}, Adv.\ Theor.\ Math.\ Phys.\ \textbf{24}
(2020) 1723--1771.

\bibitem{Gritsenko1995} V.~Gritsenko, \emph{Modular forms and moduli
spaces of abelian and K3 surfaces}, St.\ Petersburg Math.\ J.\
\textbf{6} (1995) 1179--1208.

\bibitem{GritsenkoNikulin1998} V.~Gritsenko, V.~Nikulin,
\emph{Automorphic forms and Lorentzian Kac--Moody algebras II}, Internat.\
J.\ Math.\ \textbf{9} (1998) 201--275.

\bibitem{HuangKlemm2007} M.-X.~Huang, A.~Klemm, \emph{Holomorphic
anomaly in gauge theories and matrix models}, JHEP \textbf{09} (2007) 054.

\bibitem{HKKPTVVZ2003} K.~Hori, S.~Katz, A.~Klemm, R.~Pandharipande,
R.~Thomas, C.~Vafa, R.~Vakil, E.~Zaslow, \emph{Mirror Symmetry}, Clay
Math.\ Monographs \textbf{1}, AMS (2003).

\bibitem{Kontsevich1999} M.~Kontsevich, \emph{Operads and motives in
deformation quantization}, Lett.\ Math.\ Phys.\ \textbf{48} (1999) 35--72.

\bibitem{MaulikPandharipande2006} D.~Maulik, R.~Pandharipande,
\emph{A topological view of Gromov--Witten theory}, Topology
\textbf{45} (2006) 887--918.

\bibitem{MNOP2006} D.~Maulik, N.~Nekrasov, A.~Okounkov, R.~Pandharipande,
\emph{Gromov--Witten theory and Donaldson--Thomas theory I, II}, Compos.\
Math.\ \textbf{142} (2006) 1263--1304, 1286--1304.

\bibitem{Oberdieck2017} G.~Oberdieck, \emph{Gromov--Witten invariants
of the Hilbert schemes of points of a K3 surface}, arXiv:1406.1139,
Geom.\ Topol.\ \textbf{22} (2018) 323--437.

\bibitem{OberdieckPixton2018} G.~Oberdieck, A.~Pixton,
\emph{Holomorphic anomaly equations and the Igusa cusp form
conjecture}, Invent.\ Math.\ \textbf{213} (2018) 507--587.

\bibitem{Tamarkin2007} D.~Tamarkin, \emph{Quantization of Lie bialgebras
via the formality of the operad of little disks}, GAFA \textbf{17}
(2007) 537--604.

\bibitem{YamaguchiYau2004} S.~Yamaguchi, S.-T.~Yau,
\emph{Topological string partition functions as polynomials},
JHEP \textbf{07} (2004) 047.

\end{thebibliography}

\end{document}
